We note that here we depart from the usual terminology (cf. [1]), where the name @@M0@@-ring was used for what we call a pre-@@M1@@-ring, and ``special'' @@M2@@-rings for our @@M3@@-rings; the reason is that only ``special'' @@M4@@-rings are usable in practice.

\sgaStableTarget{sga6.e05.ex.2-5}\textbf{Examples 2.5.}
\begin{enumerate}[label=\alph*)]
\item We have determined all the pre-@@M5@@-ring structures on @@M6@@; it follows from 2.4.1 that only the one defined by @@M7@@ can be a @@M8@@-ring structure. We then have @@M9@@, hence:
@@M10@@
@@M11@@
Consequently, this is indeed a @@M12@@-ring structure. Denoting by @@M13@@ the polynomial with rational coefficients @@M14@@, we therefore have:
\sgaStableTarget{sga6.e05.eq.2-5-1}@@M15@@

\item For every ring @@M16@@, the pre-@@M17@@-ring @@M18@@ is a @@M19@@-ring. First, @@M20@@. We verify the other relations by the commutativity of the diagrams:
@@M21@@
